\section{制度边界：研究输入与审阅焦点}

任何从全球机构策略翻译而来的个人框架，都会在市场制度处发生改变。交易所规则、登记结算安排、产品文件、期货公司合同、投资者适当性、数据使用权、税务和监管解释均可能随品种和日期而异。本章只列出应被核验的边界，不构成法律意见、账户资格判断、产品适当性意见或规避规则的建议。未核验即为 \qzero{}。

\begin{exhibitbox}[Exhibit 10：三类研究对象的制度核验清单]
\small
\begin{tabularx}{\exhibitboxwidth}{L{2.45cm}X L{3.9cm}}
\toprule
\textbf{对象} & \textbf{必须由当期材料核验的事项} & \textbf{本报告的硬边界} \\
\midrule
A 股 / 境内证券 & 交易规则、登记结算、T+1 或适用例外、价格限制、停复牌、风险警示、产品合同、交易成本和账户权限。 & 不把价格阈值写成必然即时的止损；不把某类 ETF 的机制外推到所有证券或基金。 \\
境内期货 & 实际合约上市状态、交易所规则、合约规格、保证金、限仓、涨跌停、交割、账户/品种权限和期货公司要求。 & 不假设某人可交易所有合约；不把连续图、理论保证金或历史合约表当作当前资格。 \\
虚拟资产相关研究 & 当期法律/监管、居民身份、场所、产品、资金、托管、数据、税务及跨境事实。 & 不构建面向中国内地居民的实施、绕行、代持、VPN、稳定币转换或境外平台使用路径。 \\
数据与研究工具 & 数据来源许可、再分发限制、抓取时间、字段定义、修订史、存储和使用范围。 & 不将无权使用、不可追溯或事后修订的数据写成可复核回测证据。 \\
\bottomrule
\end{tabularx}
\sourcenote{S11--S15。不同市场/产品的最终适用规则须以当日有效的一手材料和专业意见为准。}
\end{exhibitbox}

\section{A 股：结算、价格限制和流动性共同改变风险控制}

对普通证券的处置能力不能从一条抽象的“止损规则”推断出来。交易与登记结算规则、品种例外、风险警示、停复牌、价格限制和账户条件都会影响买入、卖出及资金可用的时间。即使某个产品存在申购赎回或其他机制，二级市场退出仍须按实际价格、数量和时点核验。故 A 股因子和事件研究必须将处置受限设为起始情景；它属于基础市场条件。

这也解释了为何本报告不输出推荐比例：一个比例若只在“可以及时卖出”的假设下安全，就无法称为经过制度约束的风险预算。研究者应该保存对应品种的当前规则版本、产品文件、交易状态和流动性证明；缺任一项，相关候选保持 \qzero{}。

\section{境内期货：存在合约不等于存在个人权限}

中国金融期货交易所的产品目录和商品期货交易所的交易规则可以帮助识别公开列示的产品与合约制度，但它们不自动构成任何个人的账户权限、适当性结论或可交易承诺。实际交易还可能涉及适当性、知识测试、资金、风险披露、期货公司合同、品种/期权限额、交易所临时通知和保证金调整等条件。研究系统只能记录“待核验”，而不能用理论合约表填补这些条件。

尤其对于 CTA 和期限结构研究，合约乘数、交割月、限仓、涨跌停、保证金和交易状态必须与日期绑定。一个今日可见的规则或合约清单不能反向证明过去某一日的条件；未来条件仍须独立核验。S11--S12、S15 的作用是指出应有的官方核验路径；个人仍须完成与自身相关的核验。

\section{虚拟资产：保持狭义、当前且不越界的表述}

中国证监会于 2026 年 2 月发布的联合通知是本备忘录在虚拟货币相关业务问题上的主要制度锚点 [S14]。本报告仅采用其狭义结论：\textbf{不把任何虚拟货币相关业务、相关金融产品交易或向中国内地居民提供的相关服务设计为本报告的实施路径。} 个人法律意见以及每一种分布式账本、数字化基础设施或经批准的现实世界资产代币化安排的法律结论，均须由独立专业审查提供。

因此，本报告的制度陈述仅覆盖 S14 所界定的虚拟货币相关业务。任何可能存在的例外、技术标签、境外主体、代理、VPN、稳定币、链上地址或平台安排，均须接受合规、账户、税务和法律审查，且不进入本报告的研究规避路径。对任何涉虚拟资产的价格/数据分析，唯一允许的报告状态是 \researchonly{} / \qzero{}。为避免歧义：\textbf{本报告不为中国大陆提供加密资产或虚拟货币相关的实施路径。}

\section{非公开信息、操纵与“看起来很聪明”的边界}

机构策略可能拥有广泛的研究、公司访问和交易对手网络，但个人研究不得将非公开重大信息、内幕消息、操纵性线索、协同交易安排或规避市场规则的操作写入策略。一个模型或社群信号如果需要信息优势、虚假深度、影响价格、隐瞒持仓、规避身份或诱导他人交易才能成立，它不属于可研究的公开市场框架。应停止记录并寻求合格的合规/法律指引。

\begin{riskbox}[规则核验的正确顺序]
先核验产品/合约和账户是否在范围内，再讨论数据和策略；先核验数据是否有权使用、当时可得，再讨论回放；先核验真实市场是否允许处置，再讨论风险阈值。把合规和市场制度放在最后，会把一个不成立的实施前提伪装成优秀模型。
\end{riskbox}
